\FloatBarrier
\newpage
\section{函数的可积性*}
\label{sec:calculus-integrability}

上一节列出了连续函数、有界单调函数和只有有限个间断点的有界函数。本节先建立共同的可积性判据，再分别研究这几类函数怎样控制取样误差。这样既能完成存在性证明，也能说明各类条件在证明中承担什么作用。

\subsection{上下和}

连续函数在每个闭子区间上取得最值，一般有界函数却未必如此。为比较全部取样值，可以使用最紧的上下界。

\begin{definition}
非空实数集合\(E\)的最小上界称为它的{\heiti{}上确界}，记为\(\sup E\)；最大下界称为它的{\heiti{}下确界}，记为\(\inf E\)。若最大值存在，它就等于上确界；上确界本身不必属于\(E\)。
\end{definition}

例如，\((0,1)\)没有最大值和最小值，但上确界为\(1\)、下确界为\(0\)。非空有界实数集合总有上下确界，这可以由第一章的单调有界原理得到：从一个非上界\(l_0\)和一个上界\(u_0\)出发，每次二分，若中点是上界就用它代替\(u_n\)，否则代替\(l_n\)。于是\(l_n\)递增、\(u_n\)递减，间距趋零，共同极限\(M\)仍是上界，而任何小于\(M\)的数最终小于某个非上界\(l_n\)，故不可能是上界。下确界同理。

设\(f\)在\([a,b]\)上有界。对分割\(P\)，定义
\begin{equation*}
 m_i=\inf_{[x_{i-1},x_i]}f,\qquad M_i=\sup_{[x_{i-1},x_i]}f,
 \qquad \omega_i=M_i-m_i.
\end{equation*}
称\(\omega_i\)为该段上的{\heiti{}振幅}。它描述全部函数值的波动范围，不要求上下界都能取到。定义{\heiti{}下和}与{\heiti{}上和}
\begin{equation*}
 L(f;P)=\sum_i m_i\Delta x_i,\qquad U(f;P)=\sum_i M_i\Delta x_i.
\end{equation*}
任意取样和满足\(L(f;P)\le S(f;P,\xi)\le U(f;P)\)，而上下和之差就是\(\sum_i\omega_i\Delta x_i\)。

\begin{figure}[!ht]
\centering
\includegraphics[width=.49\linewidth]{calculus-darboux}
\caption{蓝灰色矩形给出下和，叠加橙色部分得到上和；两者之差累积了各小段的振幅乘段长。}
\label{fig:calculus-darboux}
\end{figure}

若加点得到分割\(R\)，子段下界不低于原段下界，上界不高于原段上界，所以
\begin{equation*}
 L(f;P)\le L(f;R)\le U(f;R)\le U(f;P).
\end{equation*}
任意两个分割都可以把分点合在一起，得到共同加细分割。因此一个分割的下和不会超过另一个分割的上和。

\begin{example}
对\(f(x)=x^2\)，求\([0,1]\)上等分分割的上下和之差。再考察只在\(x=1/2\)处取\(1\)、其余位置为零的函数：哪些小段有非零振幅？说明这些小段的总贡献为何趋零。
\end{example}

\begingroup
\setlength{\parskip}{4pt}
\setlength{\abovedisplayskip}{10pt}
\setlength{\belowdisplayskip}{10pt}
\subsection{可积性判据}

\begin{theorem}[Darboux判据]\label{thm:calculus-darboux}
有界函数\(f\)在\([a,b]\)上Riemann可积，当且仅当对任意\(\varepsilon>0\)，都存在一个分割\(P\)，使
\begin{equation}\label{eq:calculus-darboux-criterion}
 U(f;P)-L(f;P)<\varepsilon.
\end{equation}
\end{theorem}

这个判据只要求找到上下和接近的分割，而积分定义要求任意足够细的分割及取样都接近同一个数。两者之间还需要一道比较。

\begin{proof}
先证必要性。若积分为\(J\)，足够细的分割上，所有取样和都与\(J\)相差小于\(\varepsilon/4\)。在每段选取接近上确界或下确界的函数值，可以让取样和分别任意接近上和或下和。因此\(U-L\le\varepsilon/2<\varepsilon\)。

再证充分性。设\(|f|\le B\)，先选定分割\(P_0\)，使上下和之差很小，记其内部分点个数为\(m\)。对另一个分割\(Q\)，加入\(P_0\)的分点得到\(R\)。除跨过这些固定分点的小段外，\(Q\)与\(R\)没有区别；受影响小段的总长度不超过\(m\|Q\|\)，每段振幅不超过\(2B\)，故
\begin{equation*}
 \begin{aligned}
 U(f;Q)-L(f;Q)&\le U(f;R)-L(f;R)+2Bm\|Q\|\\
 &\le U(f;P_0)-L(f;P_0)+2Bm\|Q\|.
 \end{aligned}
\end{equation*}
先使第一项小，再使\(\|Q\|\)足够小，就说明任意足够细的分割都有任意小的上下和之差。

对两个这样的细分割\(P,Q\)，在共同加细分割上任选取样和\(T\)。它与\(S(f;P,\xi)\)都夹在\(P\)的上下和之间，对\(Q\)也一样，所以
\begin{equation*}
 |S(f;P,\xi)-S(f;Q,\zeta)|
 \le U(f;P)-L(f;P)+U(f;Q)-L(f;Q).
\end{equation*}
右端可以任意小，且不依赖取样。

取一列最大子区间长度趋零的等分分割，在每个分割上任选和式\(S_n\)。上式表明\(\{S_n\}\)满足Cauchy准则，因而收敛于某个\(J\)。给定\(\varepsilon>0\)，先把细分程度选到任意两个和式相差小于\(\varepsilon/2\)，再取其中一个足够靠后的\(S_n\)，使\(|S_n-J|<\varepsilon/2\)。任意满足该细分要求的和式\(S\)均有
\begin{equation*}
 |S-J|\le|S-S_n|+|S_n-J|<\varepsilon.
\end{equation*}
这正是Riemann积分的定义。
\end{proof}

判据把一个极限存在性问题转为有限和的误差问题。证明中固定分点附近的误差控制，也说明为什么不能只研究一列特别选定的分割。
\endgroup

\subsection{连续函数的可积性}

现在证明定理\ref{thm:calculus-continuous-integrable}。闭区间上的连续函数一致连续。给定\(\eta>0\)，存在同一个\(\delta>0\)，使\(|u-v|<\delta\)时总有\(|f(u)-f(v)|<\eta\)。

选择\(\|P\|<\delta\)，每小段上振幅都小于\(\eta\)，于是
\begin{equation}\label{eq:calculus-upper-lower-gap}
 U(f;P)-L(f;P)=\sum_i\omega_i\Delta x_i<\eta(b-a).
\end{equation}
令\(\eta\)足够小，便满足Darboux判据，可积性得证。这里用到的不只是每个点附近函数值接近，而是所有小段同时有同一个波动上界。

若还知道\(|f'|\le K\)，则中值定理给出\(\omega_i\le K\Delta x_i\)，从而
\begin{equation*}
 U(f;P)-L(f;P)\le K\sum_i(\Delta x_i)^2\le K(b-a)\|P\|.
\end{equation*}
这提供了可计算的细分要求，但连续函数未必可导，一般证明仍依赖一致连续性。

上述估计还控制同一分割上的取样差异。任意两组取样和都在上下和之间，故其差的绝对值不超过\(U-L\)。当\(|f'|\le K\)时，对\([a,b]\)作\(n\)等分，只须\(n>K(b-a)^2/\varepsilon\)，就能使全部取样和之间的差小于\(\varepsilon\)。

\begin{exercise}
\mbox{}\par\nobreak

（1）设\(|f(u)-f(v)|\le C\sqrt{|u-v|}\)。给出上下和之差关于\(\|P\|\)的界，并证明可积性。

（2）用上题的方法处理\(f(x)=\sqrt{x}\)、\(0\le x\le1\)。为什么不能直接在整个区间上使用一个有限的导数上界？

（3）对\(f(x)=|x-1/2|\)、\(0\le x\le1\)，估计\(n\)等分时的上下和之差，并求使任意两组取样和相差小于\(10^{-3}\)的一个充分条件。
\end{exercise}

\subsection{单调函数的可积性}

设\(f\)在\([a,b]\)上单调不减。各段的下、上确界分别为\(f(x_{i-1})\)、\(f(x_i)\)，所以
\begin{equation*}
 \begin{aligned}
 U(f;P)-L(f;P)
 &=\sum_i(f(x_i)-f(x_{i-1}))\Delta x_i\\
 &\le\|P\|\sum_i(f(x_i)-f(x_{i-1}))
 =\|P\|(f(b)-f(a)).
 \end{aligned}
\end{equation*}
分割足够细时，该差趋零，故\(f\)可积。单调不增的情形应用于\(-f\)即可。这证明了定理\ref{thm:calculus-other-integrable-classes}中的第一项。

连续函数的证明使每段振幅都小；单调函数的证明则利用各段增量相加后恰为总增量。某些包含跳点的小段振幅并不随段长趋零，但它们对和式的影响仍可控制。

\begin{example}
对上一节的阶跃函数和无穷阶梯函数\(g\)，分别解释上述估计如何证明可积性。是否需要先数清全部间断点？
\end{example}

\subsection{有限个间断点的情形}

设\(f\)有界，\(|f|\le B\)，只有有限个间断点。用有限个小区间把这些点包住，使这些区间的总长度小于\(\ell\)。在剩下的有限个闭区间上，\(f\)连续，可以各自取足够细的分割，使振幅小于\(\eta\)。把所有分点合并得到\(P\)。

对包含间断点的小区间，只使用振幅不超过\(2B\)；对其余小段，使用连续性得到的振幅界。于是
\begin{equation*}
 U(f;P)-L(f;P)\le2B\ell+\eta(b-a).
\end{equation*}
先使\(\ell\)足够小，再使\(\eta\)足够小，就满足可积性判据。这证明了定理\ref{thm:calculus-other-integrable-classes}的第二项，也包括有限分段连续的有界函数。

这一证明允许间断点附近的振幅始终很大，只要求它们占据的区间总长度很小。例如，\(\sin(1/x)\)在零点附近无法用一个常值精确近似，但\([0,\ell]\)上振幅不超过\(2\)，其上下和之差的贡献不超过\(2\ell\)。

有限点的函数值也不会改变积分。若\(q\)只在\(N\)个点可能非零，且\(|q|\le B\)，任意取样和的非零贡献只能来自接触这些点的小段，其总长度不超过\(2N\|P\|\)，故
\begin{equation*}
 |S(q;P,\xi)|\le2NB\|P\|\longrightarrow0.
\end{equation*}
所以\(q\)可积且积分为零。由线性性，任意可积函数加上这样的\(q\)，积分保持不变。

\begin{exercise}
\mbox{}\par\nobreak

（1）对\(f(x)=\sin(1/x)\)、\(x>0\)，零点补零，具体说明怎样先选\([0,\ell]\)，再在\([\ell,1]\)上细分，使上下和之差小于给定误差。

（2）某函数除有限个点外等于\(x^2\)。证明它在\([0,1]\)上可积并求积分。若改为所有有理点处任意改变取值，原结论是否还必然成立？
\end{exercise}

\subsection{可积函数的运算}

可积函数的线性组合已由和式证明可积。Darboux判据还可以处理乘积与绝对值，从而补足主线使用的运算性质。

先补充区间上的限制。若\(f\)在\([a,b]\)上可积，向一个上下和接近的分割中加入\(c\)，再截取\([a,c]\)、\([c,b]\)两部分，每部分的上下和之差都不超过整个分割的差。因此\(f\)在这两个子区间上也可积，这给出了区间可加性中各项存在的依据。

设\(|f|\le A\)、\(|g|\le B\)，在任一小段内，对两点\(u,v\)，有
\begin{equation*}
 |f(u)g(u)-f(v)g(v)|
 \le A|g(u)-g(v)|+B|f(u)-f(v)|.
\end{equation*}
故乘积的振幅满足\(\omega_{fg}\le A\omega_g+B\omega_f\)。对\(f,g\)分别选择上下和接近的分割，再取共同加细分割，便得到
\begin{equation*}
 U(fg;P)-L(fg;P)
 \le A\bigl(U(g;P)-L(g;P)\bigr)+B\bigl(U(f;P)-L(f;P)\bigr).
\end{equation*}
右端可以任意小，所以\(fg\)可积。

又由\(\bigl||f(u)|-|f(v)|\bigr|\le|f(u)-f(v)|\)，知\(\omega_{|f|}\le\omega_f\)，故\(|f|\)可积。于是\(f^2\)、\(\max\{f,g\}\)、\(\min\{f,g\}\)也可积，后两者可写成\((f+g\pm|f-g|)/2\)。

这些结论说明\(\mathcal R[a,b]\)在常用运算下具有稳定的结构。连续函数还有额外性质：对连续\(f\)，\(\int f^2=0\)能推出\(f\equiv0\)；若只要求可积，单点取值的改变就使这一推断失效。

\begin{exercise}
\mbox{}\par\nobreak

（1）若\(f,g\)可积，证明\(|f-g|\)、\((f+g)^2\)可积。说明积分\(\int fg\)与\((\int f)(\int g)\)一般没有相等关系。

（2）构造不可积函数\(h\)，使\(|h|\)和\(h^2\)都可积。可考虑在有理点、无理点取不同符号。

（3）对上一节的\(D\)，直接求任意分割的上和与下和，用判据再次说明它不可积。
\end{exercise}
